
\section{MLton/MaPLe Compilation Pipeline}

\subsection{Pipeline Overview}
MLton is a whole-program optimizing compiler for Standard ML introduced in
\citet{mlton-proj}, with a parallel variant MaPLe introduced in
\citet{mpl-orig}. MaPLe expands on MLton with a variety of primitive operations
to support parallel programming (e.g. fork/join, atomics) along with runtime
support for managing pools of green threads and a parallel garbage collector.
For the purposes of this document, MLton and MPL broadly interchangeable,
differing primarily in terms of the benchmark workloads that they support:
unless explicitly noted otherwise, any statements about one compiler will also
apply to the other.

The MLton compilation pipeline passes through 8 distinct IRs
\cite{mlton-website} on the way to codegen, illustrated in
Figure~\ref{fig:mlton-ir-pipeline}:

%% http://mlton.org/IntermediateLanguage

\begin{figure}[htbp]
  \centering
  \begin{tikzpicture}[
    irnode/.style={
      rectangle, rounded corners=6pt,
      draw=colorXBorder, fill=colorX, line width=0.75pt,
      minimum height=0.75cm,
      inner xsep=8pt, inner ysep=4pt,
      font=\small\ttfamily,
      text=black!90
    },
    arr/.style={-{Stealth[length=4.5pt, width=3.5pt]}, line width=0.8pt, draw=black!55}
  ]
    \node[irnode] (ast) {AST};
    \node[irnode, right=0.4cm of ast] (coreml) {CoreML};
    \node[irnode, right=0.4cm of coreml] (xml) {XML};
    \node[irnode, right=0.4cm of xml] (sxml) {SXML};
    \node[irnode, right=0.4cm of sxml] (ssa) {SSA};
    \node[irnode, right=0.4cm of ssa] (ssa2) {SSA2};
    \node[irnode, right=0.4cm of ssa2] (rssa) {RSSA};
    \node[irnode, right=0.4cm of rssa] (mach) {Machine};

    \draw[arr] (ast) -- (coreml);
    \draw[arr] (coreml) -- (xml);
    \draw[arr] (xml) -- (sxml);
    \draw[arr] (sxml) -- (ssa);
    \draw[arr] (ssa) -- (ssa2);
    \draw[arr] (ssa2) -- (rssa);
    \draw[arr] (rssa) -- (mach);
  \end{tikzpicture}
  \caption{The MLton intermediate representation (IR) compilation pipeline.}
  \label{fig:mlton-ir-pipeline}
\end{figure}

\begin{itemize}
\item \texttt{AST}: isomorphic to the source
\item \texttt{CoreML}: removes module constructs (\texttt{structure}, \texttt{signature},
  \texttt{functor}) from \texttt{AST}
\item \texttt{XML} (``eXplicitly typed ML''): adds explicit typing to \texttt{CoreML} (but retains
  polymorphism)
\item \texttt{SXML} (``Simply typed \texttt{XML}''): removes polymorphism (but
  retains higher-order functions)
\item \texttt{SSA} (``Single Static Assignment''): removes higher-order functions,
  guarntees the usual SSA properties (as described in e.g. \citet{ssa-paper})
\item \texttt{SSA2}: optimizes representations of \texttt{ref}s and container types
\item \texttt{RSSA}: introduces explicit memory layouts
\item \texttt{Machine}: 1:1 with generated code
\end{itemize}

MLton supports a variety of codgen backends (including x86 and LLVM). In this
document, we will consider only on the C backend, since it is the only backend
supported by MaPLe.

The optimizations that we will explore target the \texttt{SSA} representation,
so in the following sections we will give a high-level overview of the mapping
from source-level SML to \texttt{SSA} (in order to capture the source-level
constructs that our optimizations target) as well as the mapping from
\texttt{SSA} to \texttt{Machine} (in order to motivate the expected performance
impact of our optimizations).

\subsection{SML $\Rightarrow$ \texttt{SSA}}
%% Regenerating code in this section:
%% 
%% Compile command:
%% ➜  lit-tests git:(f6e8544f3) /bin/rm -rf /tmp/out/*; OVERRIDE_OUTDIR=/tmp/out tools/mlton-compile.sh -verbose 3 -keep-pass '.*' -diag-pass '.*' examples/fib.sml &> /tmp/log
%%
%% Also note that the SSA IR is defined in
%% https://github.com/MLton/mlton/blob/e44f16ca8cbe1eea16657bf5bd9208f7de158436/mlton/ssa/ssa-tree.sig#L15-L245

\hl{CUT! Just say ``we'll show it this way'' and don't bother explaining what
  actually exists, it's not helpful for our purpose}

Since MLton is a whole-program compiler, it is difficult to extract small
snippets of its IRs for demonstration purposes: the trivial SML program in
\cref{fig:fib_sml} expands to a whopping 1 MB of text in the first SSA pass.
Apart from verbosity, the textual format of the IR is rather difficult to read
because it uses an unusual SML-inspired synax -- a snippet of this IR
corresponding to our example \texttt{fib} program is depicted in
\cref{fig:fib_sml_ssa}.

\begin{figure}[htbp]
  \centering
\begin{minted}{sml}
  fun fib(n) =
    if n = 0 then 1
    else n * fib(n -1)

  val _ = print (Int.toString (fib 4))
\end{minted}
  \caption{A simple \texttt{fib} program in Standard ML.}
  \label{fig:fib_sml}
\end{figure}

\begin{figure}[htbp]
  \centering
\begin{minted}{sml}
fun fib_4 (x_51048: word32): {returns = Some (word32), raises = Some (exn_27)} =
L_6736 ()
  block L_6736 ()
    val x_55402: bool = prim Word32_equal (x_51048, global_160 (*0x0:w32*))
    case x_55402 of
      true => L_7010 | false => L_1761
  block L_7010 ()
    return (global_69 (*0x1:w32*))
  block L_1761 ()
    val x_55002: word32 = prim Word32_sub (x_51048, global_69 (*0x1:w32*))
    val x_55001: bool = prim WordS32_subCheckP (x_51048, global_69 (*0x1:w32*))
    case x_55001 of
      true => L_7011 | false => L_6366
  block L_7011 ()
    raise (global_912 (*con Overflow_4*))
  block L_6366 ()
    call L_1766 (fib_4 (x_55002)) handle _ => raise
  block L_1766 (x_51053: word32)
    val x_51055: word32 = prim WordS32_mul (x_51053, x_51048)
    val x_51054: bool = prim WordS32_mulCheckP (x_51053, x_51048)
    case x_51054 of
      true => L_7011 | false => L_1767
  block L_1767 ()
    return (x_51055)
\end{minted}
  \caption{MLton's \texttt{SSA} intermediate representation for the \texttt{fib} function.}
  \label{fig:fib_sml_ssa}
\end{figure}
so we'll instead work here with a loose, Python-inspired pseudocode, and assert
that the \texttt{fib} function is lowered to \texttt{SSA} as something like the
following (where we'll adopt the simplifying convention that all functions begin
at \texttt{bb0}):
\begin{minted}{python}
  def fib(n: int): int =
    block bb0():
      isEq: boo = prim Word32_equal (n, 0)
      case isEq of true => goto bb1 | false => goto bb2

    block bb1():
      return 1

    block bb2():
      callArg: int = prim Word32_sub (n, 1)
      # call fib(n-1) and return to bb3
      call bb3 (fib (callArg))

    block bb3(callResult: int):
      funcResult: int = prim Word32_mul (callResult, n)
      return funcResult

  def main(): unit =
    block bb0():
      initVal: int = 4
      # call fib(initVal) and return to bb1
      call bb1 (fib (initVal))

    block bb1(fibResult: int):
      ...call print and exit...
\end{minted}
A few features of this representation are worth noting:

\begin{itemize}
\item Concrete operations on particular types are implemented as ``primitive
  applications,'' \texttt{prim \$PRIM\_NAME (\dots)}. \texttt{prim}s are
  eliminated at various different compilation phases, and come equipped with type
  checkers to ensure that the IR is well-formed.
\item Basic blocks carry explicit arguments. These arguments take the place of
  $\phi$ nodes in compilers for imperative languages: within each block
  \texttt{bbN} all names defined in blocks that strictly dominate \texttt{bbN}
  are available, including block arguments. To conditionally bind a value to a
  name, then, we add a conditional jump to a block that takes the desired name
  as an argument, i.e.

  \begin{minted}{python}
    block bb0():
      cond: bool = ...
      lhs: int = ...
      rhs: int = ...
      case cond of true => goto bb1 | false => goto bb2

    block bb1():
      # `lhs` is available because `bb0` dominates `bb1`
      goto bb3(lhs)

    block bb2():
      # `rhs` is available because `bb0` dominates `bb2`
      goto bb3(rhs)

    block bb3(boundValue: int):
      # All blocks dominated by bb3 now have `boundValue` available
      ...
  \end{minted}

\end{itemize}

\subsection{\texttt{SSA} $\Rightarrow$ C}
Zooming out to the most abstract level, we might end up with something like the
following for the snippet above (where we'll continue to use pseudocode in order
to elide some of the syntactic complexity, now with a C++ flavor):

\hl{TODO: This seems too long}

\begin{minted}{cpp}
  GCState gc_state;  // Global garbage collector state
  Stack stack;  // Global stack

  struct Stack {
    template <typename T>
    T* Push() {
      // Append `sizeof(T)` bytes to `data` and cast
    }

    template <typename T>
    T* Peek(int offset) {
      // Copy a `T` out of `data` and shrink
    }

    void Pop(int bytes) {
      // Drop the last `bytes` of the stack
    }

    vector<int8_t> data;  // Backing memory; dynamically-resizable
  };


  // Labels for all blocks in the program
  enum Label {
    fib_bb0,
    fib_bb1,
    ...
  };

  Label fib_chunk(Label nextBlock) {
    // Locals used across all blocks
    bool isEq;
    bool callArg;
    int funcResult;
    int n;

    // Jump table for all blocks in this chunk
  doSwitchNextBlock:
    switch (nextBlock) {
      case fib_bb0: goto L_fib_bb0;
      case fib_bb1: goto L_fib_bb1;
      ...
    }

  L_fib_bb0:
    // Read function argument from `stack`
    n = stack.Pop<int>();
    isEq = Word32_equal(n, 0);  // `prim` translated to C call
    if (isEq)
      goto l_fib_bb1;
    else
      goto l_fib_bb2;

  L_fib_bb1:  // return 1
    // Write function return value to `stack`
    *stack.Push<int>() = 1;
    // Stack frame size known at compile time
    nextBlock = stack.Peek<Label>(/*offset=*/-fibStackFrameSize);
    goto doSwitchNextBlock;

  L_fib_bb1:  // recursive call
    callArg = Word32_sub(n, 1);
    // Write return label to `stack`
    *stack.Push<Label>(fib_bb2);
    // Write function argument to `stack`
    *stack.Push<int>(callArg);
    goto fib_4;

  L_fib_bb3: // return from recursive call
    // Pop the caller result off the stack
    callResult = stack.Pop<int>();
    funcResult = Word32_mul(callResult, n);
    // Write return value to stack and return to caller
    *stack.Push<int>() = funcResult;
    nextBlock = stack.Peek<Label>(/*offset=*/-fibStackFrameSize);
    goto doSwitchNextBlock;

    ... other `fibChunk` logic ...
  }

  ... other chunks ...

  ... runtime-provided `Main` ...

\end{minted}
We'll again highlight a few notable features of the generated code:

%% TODO:
%%   1. prims are C calls
%%   2. data passes in C-temps when possible, otherwise via explicit stack
%%   3. SML funcs are collapsed into chunks

\begin{itemize}
%% \item SML control flow becomes 
\item SML functions are collapsed into ``chunks'' in the generated C code
\item Within a ``chunk,'' data passes between blocks either by explicit stack
  operations, or by ``registers'' as C-function locals.
\end{itemize}

\hl{TODO: Better transition}

We'll discuss these in more detail below

%% \paragraph{Jump }

\paragraph{Chunkification}

Rather than maintaining a mapping between SML functions and C functions, the
compiler instead attempts to implement procedure calls as either a direct
\texttt{goto} (for statically-known branch targets) or a computation to find the
appropriate entry in a jump table: this strategy makes it possible to e.g.
naturally guarantee that tail-recursive SML functions do not consume a C stack
frame per call.

For maximum performance, MLton would ideally compile the entire SML program into
a single C function: C function boundaries serve no purpose in the generated
code, and so any such boundaries are pure overhead. However, C compilation time
scales superlinearly with function size, and so this strategy can be
prohibitively expensive in practice (20+ minutes to compile a single benchmark).

As a workaround, MLton trades off improved runtime performance for decreased C
compile time by dividing the generated user program into ``chunks,'' each with
their own local jump table, local ``registers,'' and so on. Control flow between
chunks is orchestrated by returning the target block ID to the runtime, which
looks up the containing chunk in a global map, i.e.:
\begin{minted}{cpp}
  typedef enum Label {
    chunk_1_bb0,
    ...,
    chunk_2_bb0,
    ...
    doneLabel,  // Indicates program termination
  };

  Label Chunk_1(...) {
    ...local variables, jump table...

L_chunk_1_bb0:
  ...some computation...
  return chunk2_bb0;  // return to ``trampoline''
}

  Label Chunk_2(...) {
    ...local variables, jump table...

L_chunk_2_bb0:
  ...some computation...
  return chunk2_bb0;  // return to ``trampoline''
}

  typedef ChunkFn = ...; // Type of &Chunk_1, &Chunk_2, etc
  void MLton_trampoline(...) {  // Outer loop over executing chunks
    ChunkFn chunkMap[] = {
      &Chunk_1, // index == chunk_0_bb0,
      ...
      &Chunk_2, // index == chunk_1_bb0
      ...
    };
    Label nextBlock = chunk_0_bb0; // Distinguished 'start' block
    do {
      // Call into the chunk corresponding to the requested block
      nextBlock = *chunkMap[nextBlock](...);
    } while (nextBLock != doneLabel);
  }
\end{minted}


\paragraph{Stack vs ``registers''}
As we can see in the samples above, there are two possible mechanisms to pass
data between blocks in the generated C, either in ``registers,'' which are
compiled down to function-local C variables:
\begin{minted}{cpp}
  Label Chunk(...) {
    int x;
    int y;
    ...
   L_bb0:
    x = 1; // write to a "register"
    goto L_bb1;

   L_bb1:
    y = f(x); // read from a "register"
    ...
\end{minted}
or through manipulation of an explicit C-level \texttt{Stack} data structure:
\begin{minted}{cpp}
  Label Chunk(...) {
    ....
   L_bb0:
    *Stack.Push<int>() = 1; // Bump `top` pointer and initialize memory
    goto L_bb1;

   L_bb1:
    y = Stack.Peek<int>();  // Sized read from `top` pointer
    ...
\end{minted}

It is mandatory to pass data via the \texttt{Stack} in the following situations:

\begin{itemize}
\item All calls to non-tail-recursive functions: a non-tail-recursive function
  necessarily passes e.g. the return address to the callee via the stack, and so
  for the sake of simplicity, MLton chooses to pass all arguments across the
  stack in this case. \hl{TODO: CHECK}
 
\item All function calls across chunk boundaries: all generated chunks take the
  same arguments (GC state, stack state, and the target block), so there is no
  way to pass data across chunks in ``registers.''

\item All values that remain live across a GC point: for simplicity, the garbage
  collector does not examine ``registers'' when tracing live objects, hence all
  data must be spilled to the stack before the GC runs.

\end{itemize}

It is worth noting, furthermore, that passing data via C temporaries does not
\textit{guarantee} that the data remains in host machine registers: it merely
offers the C compiler's register allocator the opportunity to do so. MLton makes
no attempt to pack temporaries into available registers, and so blocks with a
large number of live-in/live-out variables will pay a performance penalty to
spill/fetch onto the (implicit, OS-level) C stack.

%% \hl{TODO}


\section{Flattening optimizations}
The following section will outline two of the existing ``flattening
optimizations'' performed by MLton. In the context of MLton, this class of
optimization was originally explored by \citet{flat-paper}. Additional
optimizations of this type were implemented by later contributors without
corresponding publications, as discussed in e.g. \citet{deep-flat-email}.

\subsection{Motivation}

\hl{CONTAINERS?}

In the source language, we have three basic species of data types (ignoring
\texttt{ref}, which is not relevant for the optimizations that we'll discuss)
\begin{itemize}
\item Primitive values: \texttt{int}, \texttt{bool}, \texttt{real}, etc
\item Tuples: a fixed-sized collection of zero or more values of a
  statically-known type, e.g. \texttt{int * int}, \texttt{real * bool * int},
  etc
\item Vectors/arrays: an immutable/mutable (respectively) variable-sized
  container for elements of a uniform, statically-known type, e.g. \texttt{int
    array}, \texttt{(int * bool) array}, etc.
\end{itemize}
In C codegen, these are translated as follows:

\hl{DIAGRAMS WOULD BE NICE}
\begin{itemize}
\item Primitive values are mapped to corresponding primitive C types, and
  reads/writes are plain C load/store operations. Primitive values can be
  allocated on the \texttt{Stack} or in local temporaries:
\begin{minted}{cpp}
  int tmp1;
  float tmp2;
  ...
  tmp1 = 42;
  tmp2 = 3.14f;
  ...
\end{minted}

\item The tuple type is mapped to a C struct type:

\begin{minted}{cpp}
  struct TupleT1 {
    struct Metadata {
      GC_header gc_header;  // Metadata tag for GC
    } metadata;
    struct Members {
      int field1;
      float field2;
      ...
    } members;
  };
\end{minted}

\hl{Add a diagram of the struct layout |header| |field1, field2,...| with arrow
  pointing to `t1`}

Unlike primitive types, tuple types are always allocated on the heap. Any
references in the generated code read/write through a pointer, e.g.
\begin{minted}{cpp}
  TupleT1::Members* t1 = nullptr; // Tuple of type T1
  float tmp;
  ...
  tmp = *bit_cast<float*>(bit_cast<uint8_t*>(t1) +
                          offsetof(TupleT1::Members, field2));
  *bit_cast<int*>(bit_cast<uint8_t*>(t1) +
                  offsetof(TupleT1::Members, field1)) = 13;
\end{minted}

\item Similarly to tuples, array types and vector types are translated to C
  struct types:

\begin{minted}{cpp}
  struct ArrTypeT1 {
    struct Metadata {
      ... GC info, length ...
    } metadata;
    struct Element {
      int el;
    } elements[];  // "Flex array" extends according to array length
  };
\end{minted}
and operations are always through pointers to data allocated on the heap:
\begin{minted}{cpp}
  ArrTypeT1::Element* arr = nullptr;  // Array of type T1
  int length;
  int tmp;
  ...
  // Element read
  tmp = bit_cast<ArrTypeT1::Element*>(bit_cast<uint8_t*>(arr) +
                                      sizeof(ArrTypeT1::Element) * 15)->el;
  // Metadata read
  length = bit_cast<ArrTypeT1::Metadata*>(bit_cast<uint8_t*>(arr) -
                                          offsetof(ArrTypeT1, Element))->length;

\end{minted}

\end{itemize}

For both the tuple and container case, nested heap-allocated data structures are
represented by \texttt{void*} members, i.e.:
\begin{minted}{cpp}
  struct TupleT2 {
    ...
    struct Members {
      void* field1;
      ...
\end{minted}
Because MLton is a whole-program compiler, all types are known statically: there
is no runtime polymorphism in the generated code. Headers carrying type
information exist only for the benefit of the GC. \hl{CUT?}


\subsection{\texttt{Flatten}}
In the source language, all functions are unary, with support for multiple
arguments implemented via tuple types, i.e. for:
\begin{minted}{sml}
  fun add(a, b) = a + b
\end{minted}
The function \texttt{add} has the type \texttt{(int * int) -> int}. This
convention is desirable in terms of simplifying language design, by removing the
need for a \hl{special construct to represent function arguments} outside of the
regular type system.

However, naively lowering \hl{this construct} to C according to the rules that
we have described above has disasterous performance characteristics: every
single function call must
\begin{enumerate}
\item Allocate a new tuple on the heap to store the function arguments
\item Copy the function arguments into the heap-allocated memory
\item Pass the address of the new tuple to the callee (on either the stack or a temporary)
\item Jump to the callee
\item Read the call arguments out of the tuple, into local variables
\item (Eventually) garbage collect the allocated tuple
\end{enumerate}
These penalties apply even in the fast intra-chunk call case. Compared to the
usual x86\_64 calling convention (where up to 14 values may pass ``for free'' in
registers \citet{amd64-abi}), we pay the price of a heap allocation per function
call and two memory accesses per argument, hurting locality and wasting memory
bandwidth.

The \texttt{Flatten} transformation of \citet{flat-paper} operates on the SSA
representation, which supports true $n$-ary functions and blocks, contrasting:
\begin{minted}{python}
  # Unary tuple-based representation
  def add_tuple(ab: (int * int)): int =
    block bb0():
      a: int = select (ab, 0)  # tuple member read
      b: int = select (ab, 1)  # tuple member read
      ret: int = prim Word32_add (a, b)
      return ret
\end{minted}
with:
\begin{minted}{python}
  # Binary "flat" representation
  def add(a: int, b: int): int =
    block bb0():
      ret: int = prim Word32_add (a, b)
      return ret
\end{minted}

A naive solution to the problem would be to interpret all SML-level functions as
$n$-ary functions over their arguments. However, this strategy can introduce
unnecessary pack/unpack operations in cases where the natural representation for
data really is a tuple, e.g.:
\begin{minted}{sml}
  val toArray: 'a array = .... (* global array value *)

  fun writeFirstElementTo (el: 'a) =
      Array.update (toArray, 0, el)

  fun readFirstElement (arr: 'a array): 'a =
    Array.sub (fromArray, 0)

  fun copyFirstElement (fromArray: 'a array) =
    writeFirstElementTo (readFirstElement fromArray)
\end{minted}
For the ``non-flat'' representation, the cost of \texttt{copyFirstElement} is
independent of the \texttt{'a} type that we instantiate the function at: even
for 1KB tuple elements, all we need to copy is a pointer to the tuple.

For the ``flat'' representation, however, the cost of \texttt{copyFirstElement}
is linear in the size of \texttt{'a} -- copying a 1KB tuple requires
``unpacking'' the entire 1KB via writes to the \texttt{Stack}, then
``re-packing'' via other 1KB worth of writes into a newly-allocated tuple.

\hl{BAD SUMMARY: DOESN'T MATCH PAPER. LET'S CIRCLE BACK AFTER DISCUSSING OUR OWN
  VERSION}

The insight of \citet{flat-paper} is to adopt a lattice-based analysis that
ensures that, for every tuple type \texttt{t}, we flatten \texttt{t} only when
we can guarantee that all uses are compatible with flattening, avoiding
pathogical cases (like the one above) that result from applying purely-local
heuristics.

This approach has the advantage that it handles more general \hl{cases of
  tuples}, including nested tuples, e.g. supporting a transformation like:
\begin{minted}{python}
  fun add_vec(input: ((int * int) * (int * int))): int =
    block bb0():
      lhs: int * int = select (input, 0)
      rhs: int * int = select (input, 1)
      lhs_x: int = select (lhs, 0)
      lhs_y: int = select (lhs, 1)
      rhs_x: int = select (rhs, 0)
      rhs_y: int = select (rhs, 1)
      ... add elements ...
\end{minted}
mapping to:
\begin{minted}{python}
  fun add_vec_flat(lhs_x: int, lhs_y: int, rhs_x: int, rhs_y: int): int =
    block bb0():
      ... add elements ...
\end{minted}
It is worth noting, however, that due to the structure of the generated C
discussed above, the savings may be less large in practice \hl{than they appear}
based on the transformation in the SSA format: a naive reading suggests that by
translating from \texttt{add\_vec} to \texttt{add\_vec\_flat}, we have saved a
total of 6 writes into temporary tuples (2 per inner tuple and 2 for the outer
tuple), eliding them naturally into register writes.

This is not true in the general case, however: if a call to
\texttt{add\_vec\_flat} crosses a C chunk boundary, then we still need to pay the
cost to write all 4 values to the \texttt{Stack}, and pay the same cost again to
pass the data on to any further flattened callers. Therefore, for nested tuple
types, it is perhaps more correct to think of flattening as a tradeoff of
bandwidth for locality (passing data through hot \texttt{Stack} memory rather
than cold heap) instead of a pure savings.

\hl{TODO: BAD DISCUSSION ABOVE; WE DON'T NECESSARILY CONSTRUCT A TEMPORARY
  TUPLE}

%% TODO:
%%   flatten: 'peels apart' tuples (in arguments)

%%     unintuitively, because of the structure of the generated C, flattening
%%     *does not* guarantee that data is passed in registers, even when the
%%     argument count is small (unlike the usual e.g. Linux calling convention)
%%
%%     in the general case it's really a tradeoff between BW and locality (the
%%     stack is always "hot" but we need to copy into it
%%
%%     chunk boundaries are not known until late so we need to make uncertain
%%     decisions

\subsection{\texttt{DeepFlatten}}
The \texttt{Flatten} transformation fixes the problem that the source language
does not support true $n$-ary functions; the \texttt{DeepFlatten} transformation
solves the problem that the source language does not support unboxed tuple types.

As discussed above, the normal MLton lowering for tuple types in containers is
\texttt{void*}, i.e. for an SML type like \texttt{(int * int) array}, we get the
C type:
\begin{minted}{cpp}
  struct ArrTypeT1 {
    struct Metadata {
      ... GC info, length ...
    } metadata;
    struct Element {
      // C-level type erasure
      void* el;
    } elements[]
  };
\end{minted}
Maniupulating elements in this kind of container corresponds to manipulating
pointers to pre-constructed tuples, stored out-of-line from the container
itself.

In terms of the local efficiency of the container operations themselves, this is
optimal: we pay a flat O(1) cost regardless of the contained element type, and
we never need to store more than a single copy of a particular tuple's data
regardness of how many containers we insert it into.

\hl{Add a diagram -- array with arrows pointing to scatterd boxes}

In terms of global program efficiency on modern CPU architectures, however, this
layout is disasterous. As an illustration, compare the following C++ snippets:
suppose we have a type:
\begin{minted}{cpp}
  struct Point {
    float x;
    float y;
  };
\end{minted}
The non-flat SML-like layout corresponds to:
\begin{minted}{cpp}
  float GetDistance(const vector<Point*>& lhs, const vector<Point*>& rhs) {
    float sum = 0.0;
    for (int i = 0; i < lhs.size(); ++i) {
      const float x = lhs[i]->x - rhs[i]->x;
      const float y = lhs[i]->y - rhs[i]->y;
      sum += x * x + y * y;
    }
    return sum;
  }
\end{minted}
Suppose that \texttt{lhs.size() = rhs.size() = n}. Then the memory bandwidth
requirement for this implementation is:
\begin{enumerate}
\item Cost to load $2n$ pointers: \texttt{$2n$ * sizeof(uintptr\_t)}. These
  values are laid out contiguously in memory, so we can assume perfect
  efficiency while loading them. \hl{REPHRASE}
\item Cost to load $2n$ \texttt{Point} objects: \texttt{$2n$ * cacheline size}.
  Since these \texttt{Point} objects are unlikely to have any particular spatial
  locality in the heap when travsered in their order in the \texttt{vector}s,
  the most likely case is that we only get \texttt{sizeof(Point) == 8} useful
  bytes out of each read from DRAM.
\end{enumerate}

Given that most modern CPU architectures have 64B cachelines and 8B pointers,
combining the two estimates above, we observe that we pay a total of 72B for
each 8B \texttt{Point} object (one 8B pointer plus one 64B cacheline), wasting a
whopping $\approx 89\%$ of the memory bandwidth that we expend on each input.

Aside from the bandwidth cost, there is also a latency cost: the load of the
pointer data (e.g. \texttt{lhs[i]->x}) cannot be pipelined with the load of the
pointer itself -- no matter how cleverly we design our out-of-order CPU, we
simply cannot issue a load to an address that we do not yet know. We thus must
wait at least a few cycles to complete a load in L1, risking a hit to our
instructions-per-cycle without heroics from the C compiler. \hl{TODO: REFERENCE}

\hl{Add a matching diagram}

In C++, the solution to this problem is to eliminate the indirection in the
source language, i.e.:
\begin{minted}{cpp}
  float GetDistanceFlat(const vector<Point>& lhs, const vector<Point>& rhs) {
    float sum = 0.0;
    for (int i = 0; i < lhs.size(); ++i) {
      const float x = lhs[i].x - rhs[i].x;
      const float y = lhs[i].y - rhs[i].y;
      sum += x * x + y * y;
    }
    return sum;
  }
\end{minted}
From the perspective of the memory access pattern of the consuming code, this
strategy is optimal: since the \texttt{Point} data is in a dense, contiguous
layout, we get 100\% efficient utilization of memory bandwidth, and we eliminate
the pipeline-hostile chain of dependent loads from the inner loop of the
program.

MLton's \texttt{DeepFlatten} optimization performs the equivalent transformation
-- for an \texttt{(int * int) array} corresponding to the \texttt{Point} type
above, we get:
\begin{minted}{cpp}
  struct ArrTypeT1 {
    struct Metadata {
      ... GC info, length ...
    } metadata;
    struct Element {
      // Flattened, inlined tuple members
      int x;
      int y;
    } elements[]
  };
\end{minted}
Unlike \texttt{Flatten}, which operates on the SSA representation (where
containers are either \texttt{'a array} or \texttt{'b array} -- both unary type
constructors), \texttt{DeepFlatten} operates on the SSA2 representation, where
we distinguish between a \texttt{((int * int) tuple) sequence} (for the non-flat
representation) vs \texttt{(int * int) sequence} (for the flat representation).

Like \texttt{Flatten}, \texttt{DeepFlatten} performs a lattice-based analysis,
flattening the types associated with particular variables only when all of their
uses are compatible with the transformation. In particular, 

\hl{Consider adding a small benchmark to confirm the napkin math}

\hl{If we want to discuss SSA vs SSA2, we need to flesh out SSA2 more. Consider
  just using SSA2 as our target IR for the whole thing. (downside: we actually
  target SSA with the optimizations later, it will raise awkward questions why
  we didn't use a different representation). Biggest conceptual difference: SSA2
  allows ``nested'' addressing into fields, i.e. t = A[i]; x = load(t, 1) -> x =
  A[i, 1]}

%% TODO:
%%   deepFlatten: 'squashes' tuples in arrays

%%   arrays are very performance sensitive for flattening because we pay the
%%   cost of an extra indirection on every access. arrays are not touched by the
%%   transformation above, which is motivated by function arguments

%%   relies on the particular structure defined above: a PackedRepresentation
%%   that squashes elements into the array. it runs on SSA2, a slight variant
%%   that gives us this exact property

%%   downsides: basically free when it works, but doesn't always work. also,
%%   since it runs after flatten, we can miss out on opportunities to do
%%   flattening that we would otherwise be able to do

%%   TODO: what does it target? when is it forbidden?

%% \hl{TODO: OUTLINE FLATTEN, DEEPFLATTEN}

\subsection{\texttt{Polyvariance}}
\hl{TODO: SHOULD PROBABLY MOVE THIS INTO AN SML -> SXML SECTION ABOVE}

Although not strictly a flattening optimization, the MLton \texttt{Polyvariance}
pass is also relevant to the optimizations we will explore here.
\texttt{Polyvariance} runs on the \texttt{SXML} IR, where the IR still has
higher order functions and implicit closure captures (all that has been removed
are the \texttt{module} constructs, which are irrelevant from the perspective of
performance). Beginning from an SML program like:
\begin{minted}{sml}
  def doOp(l: int, r: int, f: int * int -> int): int =
    f(l, r)

  def add(l: int, r: int): int =
    l + r

  def sub(l: int, r: int): int =
    l - r

  def call1(): int =
    doOp (2, 2, add)

  def call2(): int =
    doOp (2, 2, sub)
\end{minted}
we might end up with \texttt{SXML} IR like the following (using the same
Python-style pseudocode as before):
\begin{minted}{python}
  def do_op(l: int, r: int, f: int * int -> int) -> int:
    return f(l, r)

  def add(l: int, r: int) -> int:
    return l + r

  def sub(l: int, r: int) -> int:
    return l - r

  def call1() -> int:
    return do_op(2, 2, add)

  def call2() -> bool:
    return do_op(2, 2, sub)
\end{minted}
The \texttt{Polyvariance} pass duplicates function definitions at callsites
(limited by a cost heuristic that is calculated as the product of how many times
the function is called and a measure of the size of the function related to the
number of nodes in the program graph). Thus for the preceeding example, we might
end up with:
\begin{minted}{python}
  def do_op_add(l: int, r: int, f: int * int -> int) -> int:
    return f(l, r)

  def do_op_sub(l: int, r: int, f: int * int -> int) -> int:
    return f(l, r)

  def add(l: int, r: int) -> int:
    return l + r

  def sub(l: int, r: int) -> int:
    return l - r

  def call1() -> int:
    return do_op_add(2, 2, add)

  def call2() -> int:
    return do_op_sub(2, 2, sub)
\end{minted}
The benefit of this pass that is relevant to our discussion here is the impact
on closure conversion, which takes places as part of the $\texttt{SXML}
\Rightarrow \texttt{SSA}$ conversion.

While the unary-function restriction introduces \texttt{tuple} pack/unpack pairs
around all function calls in the SSA representation, \hl{closure capture}
introduces \texttt{datatype} pack/unpack pairs around calls to higher-order
functions specifically. So for the pre-\texttt{Polyvariance} snippet from above,
we might end up with the following SSA IR:
\begin{minted}{python}
  datatype lambda1 = xEnv_add | xEnv_sub
  def doOp(l: int, r: int, env: lambda1) -> int:
    block bb0():
      case env of xEnv_add => goto bb1 | xEnv_sub => goto bb2

    block bb1():
      result: int = prim Word32_add (l, r)
      return result

    block bb2():
      result: int = prim Word32_sub (l, r)
      return result
\end{minted}
The contents of each \texttt{datatype} case corresponds to the capture
environment of the function argument that it maps to, e.g. for a call like:
\begin{minted}{sml}
  fun buildAddWithBias(bias: int): int -> int = let
    fun addWithBias (l: int, r: int): int =
      return l + r + bias
  ...

  fun call3() -> int:
    doOp (2, 2, buildAddWithBias 4)
\end{minted}
we would end up with a case in the declaration like:
\begin{minted}{python}
  datatype lambda1 = xEnv_addWithBias of int | xEnv_add | xEnv_sub
\end{minted}
and a callsite like:
\begin{minted}{python}
  def call3() -> int:
    block bb0():
      biasArg: int = 5
      env: lambda_1 = xEnv_addWithBias (biasArg)
      call bbN doOp (2, 2, env)
    ...
\end{minted}
In the generated C, \texttt{datatype} objects correspond to tagged unions,
roughly:
\begin{minted}{cpp}
  struct Lambda1 {
    enum Case {
      xEnv_addWithBias,
      xEnv_add,
      xEnv_sub,
    } case;
    // Interpretation of the payload depends on the tag
    void *paylod;
  };
\end{minted}
The representation of \texttt{payload} depends on the particular
\texttt{datatype}: small types and other special cases may be allocated inline,
but the general case with a \hl{large payload} corresponds to allocating a
nested tuple, i.e. for:
\begin{minted}{sml}
  datatype env = xEnvLarge of int * int | ...
\end{minted}
the callers (in \texttt{SSA}) will look like:
\begin{minted}{python}
  l: int = ...
  r: int = ...
  t: int * int = tuple (l, r)
  e: Env = xEnvLarge t
\end{minted}
The relevance of the closure conversion transformation to flattening, then, is
that it introduces further nested tuples at callsites of higher order functions.
In the case when these nested tuples make it through to the \texttt{SSA}
representation, this indirection cannot be optimized away by \texttt{Flatten},
\texttt{DeepFlatten}, or any of the other flattening passes in MLton (since
``peeling apart'' the tuple elements is only possible when we know the value of
the tag).

The existing solution to this problem is \texttt{Polyvariance}: by duplicating
higher-order functions at their callsite, for each caller like:
\begin{minted}{python}
  def call1() -> int:
    return do_op_add(2, 2, add)

  def call2() -> int:
    return do_op_sub(2, 2, sub)
\end{minted}
the duplicated functions \texttt{do\_op\_add}, \texttt{do\_op\_sub} each end up
with a single-variant \texttt{datatype} that can easily be optimized away by
later passes (since the tag is statically known), allowing \texttt{Flatten} /
\texttt{DeepFlatten} to act on the embedded tuple as normal.

\texttt{Polyvariance} is limited, however: since it acts by duplicating
functions across all callsites, it imposes a risk of severe code bloat beyond
what is necessary just for the sake of eliminating the \texttt{datatype}. The
cost heuristic also implies the potential for a sharp performance cliff around
minor changes to the source-language program: a program that calls \texttt{map}
$N$ times might have significantly worse performance than a program that calls
\texttt{map} $N-1$ times if it pushes the \texttt{map} beyond the
\texttt{Polyvariance} threshold.


%% \hl{TODO: FINISH}

%% TODO: explain polyvariance

%%   defunctionalization runs as part of the SXML -> SSA transformation, we
%% replace higher order functions/closures with datatypes

%%   datatypes are tagged unions; the "data" portion is (maybe?) also a
%% heap-allocated tuple (give a C struct) example

%%    like how the unary-function restriction introduces tuple pack/unpack
%%  around callsites, the closure-conversion restriction of SSA introduces
%%  ConApp statements at callsites followed by case statements in callees (give
%%  a code example with a datatype + ConApp along with matching SML).

%%    these are relevant to us because they prevent tuple flattening: Flatten
%%  can't act on a type that is buried inside of a datatype (we'd need to know
%%  which case it applied to); DeepFlatten will (?) choose not to since it is an
%%  immutable member (?)

%%    the existing workaround is the polyvariance pass: according to some simple
%%  knobs, running in SXML, we (lexically) duplicate function bodies so that we
%%  end up with a single-case datatype that can be optimized away

